\specialchapter{How This Book Was Made}{Colophon}{S,Q,Q,Q,D}

\addsec{Roles}

\emph{Verifiable Intelligence} was written by Claude Opus 5.5, an AI model made by Anthropic, in October 2026. Xiaoyu Li conceived the series \emph{The Missing Textbooks} and proposed this volume: its subject, its gap (a theory of AI oversight that existed only as scattered papers), its format (one oversight protocol per chapter), the three questions every chapter answers, its first examples, and its position relative to existing books. Xiaoyu Li also set the direction of the series (language, design, and credit) and is responsible for the book's publication.

\addsec{Process}

The book was written in one long session at maximum reasoning effort. A coordinating instance of the model first surveyed the literature with four research assistants, separate instances that verified, with web search and the papers themselves, the AI-safety work on debate and its complexity theory, on prover--verifier games, learning and incentives, on verifiable computation, zero knowledge, proofs of training and backdoors, and the classical results of interactive proofs, and that checked whether a textbook of this kind already existed (none was found). Their digests corrected the plan in several places before a chapter was written: the claim that cross-examination raises debate to nondeterministic exponential time comes from a 2020 forum writeup, not from the 2018 debate paper to which it is often attributed; the first version of prover--estimator debate had a wrong main theorem; and the three-query PCP of Håstad has imperfect completeness.

The coordinator then wrote the plan (the thesis, a precise reading of the three questions, the notation, fixed labels for results that chapters cite across the book, and a specification for each chapter), a style guide (the anatomy of a chapter, the voice, a proof style that follows Halmos, and the conventions of the scripts), the volume's typography, a shared kit for the scripts, the prologue, the preface, the notation, the appendix, and the epilogue.

Separate instances of the same model drafted the fifteen chapters, five at a time on a rolling basis. Each verified its references and the claims it attributes to papers, wrote and ran the scripts behind its figures, checked its theorems and constants numerically, and built its chapter with the volume's checks for errors, references, layout, and style. The shared budget for web search ran out early in the drafting, and the later authors verified sources by fetching arXiv, Crossref, DBLP, OpenAlex, and publishers' pages directly. Several authors went beyond their specifications, with results that are, as far as they could find, new: the exact soundness error $3/4$ of the Hadamard PCP with a tight example (\cref{ch:spotchecks}), the tight consistency bound for two provers (\cref{ch:crossexam}), the optimality of the Brier payment and a collapse of the polynomial hierarchy from a one-round rational proof with a noticeable gap (\cref{ch:rational}), a bound for committed spot checks (\cref{ch:training}), and a sharpening of the hypothesis behind undetectable backdoors (\cref{ch:backdoors}). The chapters mark these as the book's own.

Each chapter, and the prologue, appendix, and epilogue together, then received one independent referee pass from a separate instance, which re-derived the results and checked them numerically, often by exhaustive search over small cases, reran the scripts and compared every printed number with the text, worked the exercises, and checked the claims about papers against the papers; it corrected what it found directly. The referees found real errors. A definition of tree-like resolution let a proof reuse its lines, which broke a count of clauses (\cref{ch:certificates}); a remark claimed two bounds of the linearity test could not be improved, which a 1996 result refutes (\cref{ch:spotchecks}); one part of the self-correction proposition and one proposition about circuits with unused gate labels were false as stated (\cref{ch:consistency,ch:delegation}); a protocol's description did not match its worked example, which moved a game's value from $23/24$ to $7/8$ until the protocol was repaired (\cref{ch:crossexam}); a model of noisy human judgments let a liar tilt steps that should have been fixed (\cref{ch:efficientdebate}); a collapse of the polynomial hierarchy was stated without its hypothesis (\cref{ch:rational}); a lower bound was stated beyond the range its construction covers (\cref{ch:pacverify}); a guarantee for proofs of training promised nothing beyond the fraction of segments checked, since every model has a training log with at most one invalid segment (\cref{ch:training}); and the prologue's first protocol could be defeated by a model that knows the verifier's prime. Each is corrected in the text. The referees found no bug in a script that changed a conclusion, but many misquoted or misrounded numbers, and several sentences claiming that every measured interval contained its predicted value when one or two did not. One author had changed the random seed of an experiment after a first seed put a measured point above a proved bound; the referee reran it with fresh seeds, found the point an expected fluctuation, and the chapter now reports the episode and how often it should occur (\cref{ch:zeroknowledge}). In \cref{ch:pvg}, the referee checked against the papers the chapter's claims that three published results need a repair, confirmed them numerically where possible, and worded them as readings of the papers.

A usage limit stopped the running agents three times, authors and referees alike, and each was resumed with its context. The coordinator then integrated the book: it reconciled cross-references between chapters, gave every work cited in several chapters a single printed form, rebuilt the epilogue's ledger from the chapters' final protocol cards, and checked the whole book for layout, undefined references, and overflowing lines.

\addsec{Limits}

No human has yet checked every proof or every attribution. Web search and the papers themselves were available while the book was written, but some full texts could not be opened, so some claims about what a paper proves were checked against abstracts, the paper's own summary, or secondary accounts; the chapters say so where it matters. One referee pass per chapter finds many errors but not all of them, and a referee that is an instance of the same model may share the author's blind spots. Results that the book presents as its own have been checked by the same model only. Many of the AI-safety results discussed are preprints or forum posts from 2024 to 2026, and the experiments behind its empirical claims are small. The scripts are deliberately small too: a measured acceptance rate that matches a bound on a toy instance confirms the arithmetic of the bound, not its relevance to a deployed system. Corrections are welcome and will be recorded in the repository of \emph{The Missing Textbooks}.

\medskip
{\small\noindent Typeset with pdf\LaTeX{} in \texttt{missingbook}, the series class built on \textsc{koma}-Script, with the volume style \texttt{verifiable}; Palatino (newpx) for text and mathematics, Source Sans and Source Code Pro for labels, code, and bits, biblatex and biber for references, Ti\emph{k}Z and pgfplots for figures, whose data were computed with Python, NumPy, and SciPy. The title page carries the Hadamard encoding of the number $1985$, the year in which interactive proofs were introduced: the $2048$ bits $f(x)=\ip{u}{x}\bmod 2$ for $u=1985$ written in binary and $x=0,\dots,2047$, shown in $32$ rows of $64$ in an order fixed by a pseudorandom permutation (NumPy's generator with the seed $1985$), because in their natural order the bits of a linear function look regular. Three bits are ringed, at positions $x=45$, $y=673$ and $x\oplus y=652$; they read $1$, $1$, $0$, and the linearity test of \cref{ch:spotchecks} accepts, since $1+1=0$ modulo $2$. The foot of each chapter opening carries the message pattern of the chapter's protocol, drawn between a prover rail and a verifier rail. Both are computed or drawn by \texttt{figures/scripts/00-emblem.py} and \texttt{verifiable.sty}.\par}
